% Oversells «Bolalar rasmli kitobi» shabloni (XeLaTeX yoki LuaLaTeX).
% Har bir bet: \rasmramka[joy=fon]{rasm}{} — rasm butun betni egallaydi
% (kesish chizigʻidan 3 mm tashqariga ham), ustida \begin{matnquti}[x=,y=,en=].
% Matn qutisini «Vizual» rejimda sichqoncha bilan surish mumkin.
\documentclass{article}
\usepackage[paperwidth=200mm,paperheight=200mm]{geometry}
\usepackage{fontspec}
\setmainfont{Noto Sans}
\usepackage{babel}
\babelprovide[import=uz-Latn, main]{uzbek}
\def\uzbekhyphenmins{22}
\usepackage{sahifa} % Oversells sahifa dizayni (sahifa.sty)

\begin{document}

% --- 1-bet: muqova ichi
\rasmramka[joy=fon]{rasmlar/quyosh.png}{}
\begin{matnquti}[x=20mm,y=30mm,en=110mm,fon=FFFFFF,shaffof=0.2]
\begin{uslub}{sarlavha}
Quyoshcha
\end{uslub}
\end{matnquti}
\clearpage

% --- 2-bet
\rasmramka[joy=fon]{rasmlar/quyosh.png}{}
\begin{matnquti}[x=20mm,y=140mm,en=160mm,fon=FFF6E0,shaffof=0.1]
\begin{uslub}{hikoya}
Har tong quyoshcha tepalik ortidan moʻralar ekan. U bogʻdagi gullarga
\matnuslubi{urgu}{«Uygʻoninglar!»} deb nur sochar ekan.
\end{uslub}
\end{matnquti}
\clearpage

% --- 3-bet
\rasmramka[joy=fon]{rasmlar/daraxt.png}{}
\begin{matnquti}[x=20mm,y=20mm,en=160mm,fon=FFFFFF,shaffof=0.15]
\begin{uslub}{hikoya}
Olma daraxti quyoshchaga rahmat aytibdi: «Sening nuring bilan
olmalarim qizarib pishdi».
\end{uslub}
\end{matnquti}
\clearpage

\end{document}
